\documentclass[10pt]{article}

\usepackage[utf8]{inputenc}

\usepackage[T1]{fontenc}

\usepackage{amsmath}

\usepackage{amsfonts}

\usepackage{amssymb}

\usepackage[version=4]{mhchem}

\usepackage{stmaryrd}

\usepackage{bbold}



\begin{document}

доставляет решение начально-краевой задачи (4), (6) и



\[

U(x, 0)=\varphi(x), \frac{\partial U}{\partial y}(x, 0)=\psi(x), 0 \leqslant x \leqslant \pi,

\]



если $a_{n}$ и $n b_{n}$ являются коәфффициентами Фурье



\[

\begin{aligned}

& a_{n}=\frac{2}{\pi} \int_{0}^{\pi} \varphi(x) \sin n x d x, \\

& n b_{n}=\frac{2}{\pi} \int_{0}^{\pi} \psi(x) \sin n x d x

\end{aligned}

\]



от достаточно гладких начальных данных $\varphi$ и $\psi$. Аналогичным образом можно получать решения начальнокраевых задач для более общих классов уравнений (1), при этом роль теории рядов Фурье, связанных с разложением (7), играет спектральная теория линейных операторов.



Ф. м. тесно связан со специальными функциями, к-рые являются решениями уравнениӥ (3) при $m=$ $=n=1$ для частных случаев операторов $M$ и $N$, многие из этих функций первоначально возникли таким сіособом. Напр., при применении этого метода к уравнению Гельмгольца



\[

u_{x x}+u_{y y}+u=0,

\]



записанного в полярных координатах в виде (1) с



\[

M=r^{2} \frac{\partial}{\partial r^{2}}+r \frac{\partial}{\partial r}+r^{2}, N=-\frac{\partial}{\partial \theta^{2}},

\]



первое уравнение в (3) является уравнением Бесселя.



Одно и то же диффференциальное уравнение имеет, вообще говоря, целое семейство систем координат, в к-рых оно допускает разделение переменных, т. е. приводится к виду (1). Задача разыскания таких систем координат тесно связана с групповыми свойствами диффференциальных уравнений. Применение методов теории групп Ли позволяет описать все решения с разделенными переменными многих классич. уравнений математич. физики (Лапласа, Гельмгольца, IÏёдингера, волнового уравнения и др.). На этом пути получается также целый ряд соотношений из теории специальных функций.



Метод разделения переменных был предложен для решения волнового уравнения Ж. Д'Аламбером (Ј. D'Alembert, 1749), с достаточной полнотой метод был развит в нач. 19 в. Ж. Фурье (J. Fourier) и в полной общности сформулирован М. В. Остроградским в 1828.



Лum. . [1] Б и ц а Д з е А. В., Уравнения математической

М., 1976, [2] М и л л е р У., Симметрия и разделение физики, М., 1976, [2] М и л л е р У., Симметрия и разделение

переменных, пер. с англ., М., 1981.

А. П. Солдатов.



ФУРЬЕ ПОКАЗАТЕЛИ П о ч т И п е р и о д и ч ес к о й ф у н ц и и - действительные числа $\lambda_{n}$ в Фурье ряде, соответствующем данной почти периодич. функции $f(x)$ :



\[

f(x) \sim \sum_{n} a_{n} e^{i \lambda_{n} x}

\]



где $a_{n}-\Phi$ Фье коәффициентьи функции $f(x)$. Множество показателей Фурье функции $f(x)$ наз. с п е кт р о м этой функции. В отличие от периодич. случая спектр почти периодич. функции может иметь предельные точки на конечном расстоянии и даже быть всюду плотным. Поэтому поведение ряда Фурье потти периодич. функции существенно зависит от арифмметич. структуры ее спектра.



E. А. Бредихина.



ФУ РБЕ ПРЕОБРАЗОВАНИЕ, одно из интегральных преобразований, - линейный оператор $F$, действующий в пространстве, әлементами к-рого являются функции $f(x)$ от $n$ действительных переменных. Минимальной областью определения $F$ считается совокупность $D=C_{0}^{\infty}$ бесконечно дифференцируемых финитных функций $\varphi$. Для таких функций



\[

(F \varphi)(x)=(2 \pi)^{-n / 2} \int_{\mathbb{R}^{n}} \varphi(\xi) e^{-i x \xi} d \xi

\]



В нек-ром смысле наиболее естественной областью определения $F$ является совокупность $S$ бесконечно дифференцируемых функций $\varphi(x)$, исчезающих на бесконечности вместе со своими производными быстрее любой степени $|x|$. Формула (1) сохраняется для $\varphi \in S$ и при этом $(F \varphi)(x) \equiv \psi(x) \in S$. Более того, $F$ осуществляет изоморфизм $S$ на себя, обратное отображение $F^{-1}$ — обращение Ф. п., о б р а т н о е п р е о б р а 3 о в ан и е Ф у р ье, - задается формулой:



\[

\varphi(x)=\left(F^{-1} \psi\right)(x)=(2 \pi)^{-n / 2} \int_{\mathbb{R}^{n}} \psi(x) e^{i x \xi} d \xi .

\]



Формула (1) еще действует в пространстве $L_{1}(\mathbb{R} n)$ суммируемых функций. Дальнейшее расширение области определения оператора $F$ требует обобщения формулы (1). В классич. анализе такие обобщения строятся для локально суммируемых функций с теми или иными ограничениями на их поведение при $|x| \rightarrow \infty$ (см. Фурье интеграл). В теории обобщенных функций определение оператора $F$ освобождено от многих требований классич. анализа.



Основные задачи, связанные с изучением Ф. п. $F$ : исследование области определения $\Phi$ оператора $F$ и области его значений $F \Phi=\Psi$, свойства отображения $F: \Phi \rightarrow \Psi$ (в частности, условия существования обратного оператора $F^{-1}$ и его выражение). Формула обращения Ф. п. весьма проста:



\[

F^{-1}[g(x)]=F[g(-x)]

\]



Под действием Ф. п. линейные операторы в исходном иространстве, инвариантные относительно сдвига, переходят в пространстве образов в операторы умножения (при нек-рых условиях). В частности, свертка функций $f$ и $g$ переходит в произведение функций $F f$ и $F g$ :



\[

F(f * g)=F f \cdot F g,

\]



дифференцирование порождает умножение на независимую переменную:



\[

F\left(D \alpha_{f}\right)=(i x)^{\alpha} F f .

\]



В пространствах $L_{p}\left(\mathbb{R}^{n}\right), 1 \leqslant p \leqslant 2$, оператор $F$ определен формулой (1) на множестве $D_{F}=L_{1} \bigcap L_{p}\left(\mathbb{R}^{n}\right)$ и является ограниченным оператором из $L_{p}(\mathbb{R} n)$ в $L_{q}(\mathbb{R} n), \frac{1}{p}+\frac{1}{a}=1$ :



\[

\begin{aligned}

& \left\{(2 \pi)^{-n / 2} \int_{\mathbb{R}^{n}}|(F f)(x)|^{q} d x\right\}^{1 / q} \leqslant \\

& \leqslant\left\{(2 \pi)^{-n / 2} \int_{\mathbb{R}^{n}}|f(x)| p d x\right\}^{1 / p}

\end{aligned}

\]



(нер авенство Х у у дор фа-Юнга). ПІо непрерывности $F$ допускает продолжение на все пространство $L_{p}(\mathbb{R} n)$, к-рое (для $1<p \leqslant 2$ ) дается формулой



\[

(F f)(x)=\lim _{R \rightarrow \infty}^{q}(2 \pi)^{-n / 2} \int_{|\xi|<R} f(\xi) e^{-i \xi x} d \xi=\tilde{f}(x),

\]



где сходимость понимается по норме пространства $L_{q}\left(\mathbb{R}^{n}\right)$. Если $p \neq 2$, образ пространства $L_{p}$ под действием оператора $F$ не совпадает с $L_{q}$, т. е. вложение $F L_{p} \subset L_{q}$ строгое при $1 \leqslant p<2$ (случай $p=2$ см. в статье Планшереля теорема). Обратный оператор $F^{-1}$ определен на $F L_{p}$ формулой



\[

\begin{gathered}

(F-1 \tilde{f})(x)=\lim _{R \rightarrow \infty}^{p}(2 \pi)^{-n / 2} \int_{|\xi|<R} \tilde{f}(\xi) e^{i \xi x} d \xi, \\

1<p \leqslant 2 .

\end{gathered}

\]



Задача о распространении (D. ІІ. на возможно пирокий класс функций постоянно возникает в анализе и его приложениях. См., напр., Фурье преобразование приложениях. См., напр.,

о б о бщ е н н о й ф у н к ц и.





\end{document}